\section{Sheafification}
\label{section-sheafification}

\noindent
In this section we explain how to get the sheafification
of a presheaf on a topological space. We will use stalks
to describe the sheafification in this case. This is different
from the general procedure described in Sites, Section
\ref{sites-section-sheafification}, and perhaps somewhat
easier to understand.

\medskip\noindent
The basic construction is the following. Let $\mathcal{F}$ be a presheaf
of sets on a topological space $X$.
For every open $U \subset X$ we define
$$
\mathcal{F}^{\#}(U)
=
\{
(s_u) \in \prod\nolimits_{u \in U} \mathcal{F}_u
\text{ such that }(*)
\}
$$
where $(*)$ is the property:
\begin{itemize}
\item[(*)] For every $u \in U$, there exists an open neighbourhood
$u \in V \subset U$, and a section $\sigma \in \mathcal{F}(V)$
such that for all $v \in V$ we have $s_v = (V, \sigma)$
in $\mathcal{F}_v$.
\end{itemize}
Note that $(*)$ is a condition for each $u \in U$,
and that given $u \in U$ the truth of this condition
depends only on the values $s_v$ for $v$ in any open neighbourhood
of $u$. Thus it is clear that,
if $V \subset U \subset X$ are open, the projection maps
$$
\prod\nolimits_{u \in U} \mathcal{F}_u
\longrightarrow
\prod\nolimits_{v \in V} \mathcal{F}_v
$$
maps elements of $\mathcal{F}^{\#}(U)$ into $\mathcal{F}^{\#}(V)$.
Using these maps as the restriction mappings, we turn $\mathcal{F}^\#$
into a presheaf of sets on $X$.

\medskip\noindent
Furthermore, the map $\mathcal{F}(U) \to \prod_{u \in U} \mathcal{F}_u$
described in Section \ref{section-stalks} clearly has image
in $\mathcal{F}^{\#}(U)$. In addition, if $V \subset U \subset X$ are
open then we have the following commutative diagram
$$
\xymatrix{
\mathcal{F}(U) \ar[r] \ar[d] &
\mathcal{F}^{\#}(U) \ar[r] \ar[d] &
\prod_{u\in U} \mathcal{F}_u \ar[d] \\
\mathcal{F}(V) \ar[r] &
\mathcal{F}^{\#}(V) \ar[r] &
\prod_{v\in V} \mathcal{F}_v
}
$$
where the vertical maps are induced from the
restriction mappings. Thus we see that
there is a canonical morphism of presheaves
$\mathcal{F} \to \mathcal{F}^{\#}$.

\medskip\noindent
In Example \ref{example-sheaf-product-pointwise} we saw
that the rule $\Pi(\mathcal{F}) : U \mapsto \prod_{u\in U} \mathcal{F}_u$
is a sheaf, with obvious restriction mappings. And by construction
$\mathcal{F}^{\#}$ is a subpresheaf of this. In other words, we
have morphisms of presheaves
$$
\mathcal{F} \to \mathcal{F}^\# \to \Pi(\mathcal{F}).
$$
In addition the rule that associates to $\mathcal{F}$
the sequence above is clearly functorial in the presheaf $\mathcal{F}$.
This notation will be
used in the proofs of the lemmas below.

\begin{lemma}
\label{lemma-sheafification-sheaf}
The presheaf $\mathcal{F}^{\#}$ is a sheaf.
\end{lemma}

\begin{proof}
It is probably better for the reader to find their own explanation
of this than to read the proof here. In fact the lemma is true
for the same reason as why the presheaf of continuous
function is a sheaf, see Example \ref{example-basic-continuous-maps}
(and this analogy can be made precise using the ``espace \'etal\'e'').

\medskip\noindent
Anyway, let $U = \bigcup U_i$ be an open covering.
Suppose that $s_i = (s_{i, u})_{u \in U_i} \in \mathcal{F}^{\#}(U_i)$
such that $s_i$ and $s_j$ agree over $U_i \cap U_j$.
Because $\Pi(\mathcal{F})$ is a sheaf,
we find an element $s = (s_u)_{u\in U}$ in $\prod_{u\in U} \mathcal{F}_u$
restricting to $s_i$ on $U_i$. We have to check property $(*)$.
Pick $u \in U$. Then $u \in U_i$ for some $i$. Hence by $(*)$ for $s_i$,
there exists a $V$ open, $u \in V \subset U_i$
and a $\sigma \in \mathcal{F}(V)$
such that $s_{i, v} = (V, \sigma)$ in $\mathcal{F}_v$
for all $v \in V$. Since $s_{i, v} = s_v$ we get $(*)$ for $s$.
\end{proof}

\begin{lemma}
\label{lemma-stalk-sheafification}
Let $X$ be a topological space.
Let $\mathcal{F}$ be a presheaf of sets on $X$.
Let $x \in X$. Then $\mathcal{F}_x = \mathcal{F}^\#_x$.
\end{lemma}

\begin{proof}
The map $\mathcal{F}_x \to \mathcal{F}^\#_x$
is injective, since already the map
$\mathcal{F}_x \to \Pi(\mathcal{F})_x$ is injective.
Namely, there is a canonical map $\Pi(\mathcal{F})_x \to \mathcal{F}_x$
which is a left inverse to the map $\mathcal{F}_x \to \Pi(\mathcal{F})_x$,
see Example \ref{example-sheaf-product-pointwise-stalk}.
To show that it is surjective, suppose that
$\overline{s} \in \mathcal{F}^\#_x$.
We can find an open neighbourhood $U$ of $x$ such that
$\overline{s}$ is the equivalence class of $(U, s)$
with $s \in \mathcal{F}^\#(U)$.
By definition, this means there exists an open neighbourhood
$V \subset U$ of $x$ and a section $\sigma \in \mathcal{F}(V)$
such that $s|_V$ is the image of $\sigma$ in $\Pi(\mathcal{F})(V)$.
Clearly the class of $(V, \sigma)$ defines an element of
$\mathcal{F}_x$ mapping to $\overline{s}$.
\end{proof}

\begin{lemma}
\label{lemma-sheafify-universal}
Let $\mathcal{F}$ be a presheaf of sets on $X$.
Any map $\mathcal{F} \to \mathcal{G}$ into a sheaf of sets
factors uniquely as
$\mathcal{F} \to \mathcal{F}^\# \to \mathcal{G}$.
\end{lemma}

\begin{proof}
Clearly, there is a commutative diagram
$$
\xymatrix{
\mathcal{F} \ar[r] \ar[d] &
\mathcal{F}^\# \ar[r] \ar[d] &
\Pi(\mathcal{F}) \ar[d] \\
\mathcal{G} \ar[r] &
\mathcal{G}^\# \ar[r] &
\Pi(\mathcal{G}) \\
}
$$
So it suffices to prove that $\mathcal{G} = \mathcal{G}^\#$.
To see this it suffices to prove, for every point $x \in X$ the
map $\mathcal{G}_x \to \mathcal{G}^\#_x$ is bijective, by
Lemma \ref{lemma-points-exactness}. And this is
Lemma \ref{lemma-stalk-sheafification} above.
\end{proof}

\noindent
This lemma really says that there is an adjoint pair
of functors: $i : \Sh(X) \to \textit{PSh}(X)$
(inclusion) and $\# : \textit{PSh}(X) \to \Sh(X)$
(sheafification). The formula is that
$$
\Mor_{\textit{PSh}(X)}(\mathcal{F}, i(\mathcal{G}))
=
\Mor_{\Sh(X)}(\mathcal{F}^\#, \mathcal{G})
$$
which says that sheafification is a left adjoint of
the inclusion functor. See Categories, Section
\ref{categories-section-adjoint}.

\begin{example}
\label{example-sheafify-constant}
See Example \ref{example-stalk-constant-presheaf} for notation.
The map $A_p \to \underline{A}$ induces a map
$A_p^\# \to \underline{A}$. It is easy to see that this
is an isomorphism. In words: The sheafification
of the constant presheaf with value $A$ is the
constant sheaf with value $A$.
\end{example}

\begin{lemma}
\label{lemma-separated-presheaf-into-sheaf}
Let $X$ be a topological space.
A presheaf $\mathcal{F}$ is separated (see
Definition \ref{definition-separated}) if and only if
the canonical map $\mathcal{F} \to \mathcal{F}^\#$ is injective.
\end{lemma}

\begin{proof}
This is clear from the construction of $\mathcal{F}^\#$ in this
section.
\end{proof}

\begin{lemma}
\label{lemma-sheafify-epi-mono}
Let $X$ be a topological space. The sheafification of a surjective
(resp.\ injective) morphism of presheaves of sets is surjective
(resp.\ injective).
\end{lemma}

\begin{proof}
Omitted.
\end{proof}
